\subsection{彼得--外尔定理}

回忆我们以 \(\Theta(\mathrm{G})\) 表示 G 的有限维酉表示的矩阵系数空间。空间 \(\Theta(\mathrm{G})\) 是 \(\mathcal{C}(\mathrm{G})\) 的一个在复共轭下稳定的含单位元子代数（V, p. 386, 命题 5）。

对 \(\pi \in \widehat{\mathrm{G}}\)，令 \(\varrho_{\mathrm{G}}(\pi)\) 为 \(\mathcal{C}(\mathrm{G})\) 中由 \(\pi\) 的矩阵系数生成的子空间。将其与 \(\mathrm{L}^2 (\mathrm{G})\) 的一个子空间等同。

空间 \(\Theta(G)\) 等于各空间 \(\varrho_{G}(\pi)\) 的和，其中 \(\pi\in\widehat{G}\)。事实上，\(\Theta(G)\) 的每个元都是 G 的不可约表示的矩阵系数之和，因为由 V, p. 457, 命题 1，G 的有限维表示是半单的。此外，该和是直和，因为由 V, p. 458, 推论 2，各空间 \(\varrho_{G}(\pi)\) 两两正交。

命题 4. — 空间 \(\Theta(\mathrm{G})\) 在 \(\mathcal{C}(\mathrm{G})\) 和 \(\mathrm{L}^2(\mathrm{G})\) 中都稠密。

含单位元子代数 \(\Theta(\mathrm{G})\) 的 \(\mathcal{C}(\mathrm{G})\) 在共轭下稳定。由 V, p. 457, 命题 2，它与子代数 \(\Upsilon(\mathrm{G})\) 重合，故能分离 G 的点（V, p. 454, 定理 4 的推论）。因此，由 TG, X, p. 40, 推论 2，它在 \(\mathcal{C}(\mathrm{G})\) 中稠密，从而在 \(\mathrm{L}^2(\mathrm{G})\) 中也稠密（参见 INT, IV, p. 155, §4, no. 7, 命题 13）。

回忆 G 的双正则酉表示是表示 \(\varrho_{G}\)，它从 G × G 到 \(\mathrm{L}^{2}(\mathrm{G})\)，满足 \((g_{1}, g_{2}) \mapsto \gamma_{\mathrm{G}}(g_{1}) \delta_{\mathrm{G}}(g_{2})\)。

命题 5. --- 设 \(\pi \in \widehat{\mathbf{G}}\)。空间 \(\varrho_{\mathrm{G}}(\pi)\) 是 \(\varrho_{\mathrm{G}}\) 的一个子表示，与 \(\overline{\pi} \boxtimes \pi\) 同构。特别地，它是 \(\mathrm{G} \times \mathrm{G}\) 的不可约表示。

由 V, p. 422, 命题 7（应用 \(Z = \{e\}\)），空间 \(\varrho_{G}(\pi)\) 是 \(\varrho_{G}\) 的子表示，且从 \(\overline{\mathrm{E}}_{\pi} \otimes \mathrm{E}_{\pi}\) 到 \(L^2(G)\)、将 \(x \otimes y\) 送到矩阵系数 \(g \mapsto \langle x | \pi(g)y \rangle\) 的线性映射，是一个 \((G \times G)\)-同构，从 \(\overline{\pi} \boxtimes \pi\) 到 \(\varrho_{G}(\pi)\)。因此子表示 \(\varrho_{G}(\pi)\) 是不可约的（V, p. 461, 命题 3）。

定理 1（彼得--外尔）。--- 设 G 是紧拓扑群。G 的双正则表示 \(\varrho_{\mathrm{G}}\) 是各子表示 \(\varrho_{\mathrm{G}}(\pi)\)（其中 \(\pi \in \widehat{\mathrm{G}}\)）的希尔伯特直和。

各空间 \(\varrho_{\mathrm{G}}(\pi)\) 两两正交；它们是 G 的双正则表示的不可约子表示（命题 5），并且两两不同构（V, p. 461, 命题 3）。

定理即由以下事实得出：空间 \(\Theta(\mathrm{G})\) 是各空间 \(\varrho_{\mathrm{G}}(\pi)\) 的和，其中 \(\pi\) 遍历 \(\widehat{\mathrm{G}}\)，且该和在 \(\mathrm{L}^2(\mathrm{G})\) 中稠密（命题 4）。

推论 1. --- 对每个 \(\pi\in\widehat{G}\)，子表示 \(\varrho_{\mathrm{G}}(\pi)\) 是 \((\overline{\pi}\boxtimes\pi)\)-同型分量的 \(\varrho_{\mathrm{G}}\)。

令 F 为 \((\overline{\pi} \boxtimes \pi)\)-同型分量的 \(\varrho_{G}\)。空间 F 包含 \(\varrho_{\mathrm{G}}(\pi)\)（命题 5）。此外，对每个 \(\tau \in \widehat{G}\)（不同于 \(\pi\)），F 与 \(\varrho_{\mathrm{G}}(\tau)\) 的交为零（V, p. 461, 命题 3）。应用定理可推出 \(\mathrm{F} = \varrho_{\mathrm{G}}(\pi)\)。

推论 2. --- 设 \(\pi_1\) 和 \(\pi_2\) 是不同构的 G 的不可约表示。\((\overline{\pi}_1 \boxtimes \pi_2)\)-同型分量的 \(\varrho_G\) 为零。

推论 3. --- G 的右（分别为左）正则表示同构于希尔伯特直和

\[
\bigoplus_ {\pi \in \widehat {\mathrm{G}}} \pi^ {\dim (\pi)}.
\]

由定理及命题 5，G 的双正则表示限制到子群 \(\mathrm{H} = \{e\} \times \mathrm{G}\) 后，同构于各表示 \(\overline{\pi} \boxtimes \pi\) 在 H 上的限制的希尔伯特直和，其中 \(\pi \in \widehat{\mathrm{G}}\)。这些表示同构于 \(\dim(\pi)\) 个 \(\pi\) 的直和（V, p. 384, 引理 8）。这给出右正则表示的结果，左正则表示的情形类似。

推论 4. --- 设 \(\pi\in\widehat{G}\)。\(\pi\)-同型分量的 \(\delta_{G}\)（分别为 \(\overline{\pi}\)-同型分量的 \(\gamma_{G}\)）等于 \(\varrho_{G}(\pi)\)。

论证与前一推论类似，只需将 \(\varrho_{\mathrm{G}}(\pi)\) 看作 \(\delta_{\mathrm{G}}\)（分别为 \(\gamma_{\mathrm{G}}\)）的子表示。

注。--- 1) 若 G 有限，这些陈述对应于 A, VIII, p. 398 的注记中的结果。

\begin{enumerate}
\def\labelenumi{\arabic{enumi})}
\setcounter{enumi}{1}
\item
  假设 G 是交换的。集合 \(\widehat{G}\) 与 G 的对偶群相等同（V, p. 393, 注记）。对每个 \(\chi \in \widehat{G}\)，空间 \(\varrho_{\mathrm{G}}(\chi)\) 是 \(\mathrm{L}^2 (\mathrm{G})\) 中由 \(\chi\) 生成的一维子空间。因此，G 的定理 1 等价于 II, p. 215 的定理 1 的推论。
\item
  若存在整数 \(n \geqslant 0\)，使得 G 是 \(\mathbf{GL}(\mathbf{C}^{n})\) 的紧子群，则 G 在 \(\mathbf{GL}(\mathbf{C}^{n})\) 中的恒等表示足以分离 G 的点，随后可以直接证明命题 4，再证明彼得–外尔定理，而无需援引盖尔范德–赖科夫定理。
\item
  彼得–外尔定理特别蕴含，从 G 到 \(\prod_{\pi\in\widehat{G}}\mathbf{U}(\mathrm{E}_{\pi})\) 的自然连续同态是单射。
\end{enumerate}

